% =====================================================================
% CAPITULO 1 - INTRODUCCION
% =====================================================================
% Los objetivos especificos de 1.3.2 estan numerados y son verificables a
% proposito: la seccion 6.1 los cierra uno a uno con la evidencia concreta.
% Si se toca uno aqui, hay que tocar su cierre alli.
%
% ESTILO (31-ago): sin negritas dentro del texto corrido, sin incisos entre
% rayas salvo cuando la frase los pida de verdad, y sin florituras.
% =====================================================================

\section{Contexto y motivación}
\label{sec:contexto}

Un sistema de detección de intrusiones en red analiza el tráfico en busca de
actividad maliciosa. Durante décadas lo hizo mediante firmas: reglas que
describen patrones de bytes asociados a amenazas ya conocidas
\cite{Swarnkar2023}. El método es preciso y su coste de ejecución es bajo, pero
tiene un límite que el mantenimiento no resuelve, porque solo detecta lo que
alguien ha analizado y catalogado previamente.

Ese límite se ha agravado por tres motivos. El tráfico circula hoy
mayoritariamente cifrado, así que el contenido que la firma inspeccionaría
resulta indistinguible del ruido. El malware es polimórfico y altera su
apariencia entre infecciones sin cambiar su función. Y las herramientas de
ataque se han vuelto accesibles, lo que multiplica el número de variantes en
circulación.

Por eso el campo ha girado hacia el aprendizaje automático, que induce el
criterio de detección a partir de los datos en lugar de recibirlo escrito. La
ventaja esperada es la capacidad de generalizar a variantes no vistas. El coste
son los falsos positivos, la opacidad del modelo y la dependencia de datos
etiquetados de calidad. Este trabajo se ocupa sobre todo de esto último, que
resultó ser el problema más costoso.

\section{El problema: medir la detección es más difícil que hacerla}
\label{sec:problema}

El planteamiento inicial de este trabajo era el habitual en la literatura de
detección híbrida y se puede enunciar como una tensión entre dos fuentes de
información.

Por un lado están las características de flujo: cuántos paquetes componen la
conexión, cuántos bytes transporta y cuánto dura. Están siempre disponibles,
funcionan igual con tráfico cifrado y son ciegas al contenido. Por otro está la
carga útil, es decir, los bytes que la conexión transporta. Contienen el ataque
cuando el ataque está en los datos, y no aportan nada cuando el tráfico va
cifrado o cuando el ataque no llega a intercambiar contenido.

La hipótesis de partida era que cada tipo de ataque escondería su firma en una
de esas dos fuentes, de manera que ninguna bastaría por sí sola y el enfoque
híbrido quedaría justificado por esa vía. Para comprobarlo se seleccionaron seis
días del conjunto CSE-CIC-IDS2018 que cubren cuatro regímenes distintos: fuerza
bruta cifrada, ataques web en claro, denegación de servicio volumétrica y una
botnet con actividad periódica.

Medida con cuidado, esa hipótesis no se sostiene, y el motivo es el hallazgo
central de esta memoria. Dentro de un mismo día, cualquiera de las
representaciones detecta cualquiera de los ataques con un macro-F1 por encima de
$0.97$, de modo que la comparación está saturada y no permite distinguirlas.
Un resultado alto obtenido así no informa sobre la capacidad de detección.

El problema que este trabajo acaba abordando está, por tanto, un nivel más
abajo: cómo se mide si un detector ha aprendido lo que se cree que ha aprendido.
La respuesta que propone es que hay que perturbar las condiciones de evaluación,
dejando que el atacante camufle sus bytes, cambiando de día o cambiando de
conjunto de datos, porque es ahí donde las representaciones dejan de ser
equivalentes.

La tesis del trabajo sobrevive, pero apoyada en otra evidencia: ninguna vista
basta por sí sola, y quienes lo demuestran son la prueba de evasión, la
transferencia entre dominios y los ataques que no dejan contenido que analizar.

\section{Objetivos}
\label{sec:objetivos}

\subsection{Objetivo general}

Evaluar, sobre datos reales de red y con un protocolo experimental explícito, en
qué condiciones el análisis de la carga útil aporta información que las
características de flujo no capturan, y en cuáles no; y determinar qué
representación resiste cuando cambian las condiciones de evaluación.

\subsection{Objetivos específicos}

Se enuncian de forma verificable, para poder cerrarlos uno a uno en la
sección~\ref{sec:objetivos-cumplidos}.

\begin{enumerate}
  \item[\textbf{O1}] Construir un \emph{pipeline} reproducible que convierta
        capturas de red en bruto en un conjunto de ejemplos etiquetados, con la
        conexión TCP como unidad y varias representaciones por conexión.
  \item[\textbf{O2}] Establecer una \emph{ground-truth} verificada sobre las
        capturas, en lugar de heredarla de la documentación del conjunto de
        datos, y documentar las discrepancias que aparezcan.
  \item[\textbf{O3}] Comparar las representaciones con un único protocolo, es
        decir, con la misma métrica, el mismo balanceo y el mismo tratamiento
        del tráfico benigno, sobre varios regímenes de ataque, y establecer si
        la comparación discrimina entre ellas.
  \item[\textbf{O4}] Determinar por qué acierta cada representación, y no solo
        cuánto, mediante interpretabilidad y pruebas de evasión activas.
  \item[\textbf{O5}] Medir la generalización entre días del mismo régimen y
        entre conjuntos de datos distintos, e identificar qué componente del
        modelo falla cuando falla.
  \item[\textbf{O6}] Contrastar arquitecturas del estado del arte con las
        propias a presupuesto de parámetros comparable, para determinar si el
        límite de los resultados está en la capacidad del modelo.
\end{enumerate}

\section{Aportaciones de este trabajo}
\label{sec:aportaciones}

Hay que separar lo que se toma de la literatura de lo que este trabajo añade,
porque la aportación no es una arquitectura nueva.

\subsection*{Lo que se toma del estado del arte}

Las representaciones de la carga útil siguen la tradición de PAYL y Anagram
\cite{Wang2004, Wang2006}. El modelo convolucional sobre bytes sigue la línea de
\emph{Deep Packet} \cite{Lotfollahi2020}. El \emph{transformer} sobre bytes y la
fusión por atención cruzada se implementan siguiendo sus referencias
respectivas. El escalado de Platt \cite{Platt1999} y el error de calibración
esperado \cite{Guo2017} son métodos estándar. La crítica a los conjuntos de
datos públicos parte de Engelen \emph{et al.} \cite{Engelen2021}. Todo ello se
señala en el punto donde se usa.

\subsection*{Lo que aporta este trabajo}

Son cinco resultados, agrupados por lo que cada uno permite afirmar.

El primero es una \emph{ground-truth} verificada sobre las capturas para siete
días, que corrige la documentación oficial en tres de ellos: la dirección del
atacante SSH, el puerto del ataque lento y la continuidad de la actividad de la
botnet (sección~\ref{sec:ground-truth}).

El segundo es la constatación de que la evaluación intra-día está saturada y no
discrimina entre representaciones (sección~\ref{sec:saturacion}). Se propuso una
causa, que un único host atacante por día permitiera memorizar su identidad, y
se contrastó y refutó con un experimento de exclusión de atacantes que incluye
control positivo (sección~\ref{sec:por-que-satura}).

El tercero es la medición de que el éxito del análisis de contenido suele ser
una huella de la herramienta y no del ataque, demostrada por interpretabilidad y
por evasión (secciones~\ref{sec:saliency} y~\ref{sec:evasion}).

El cuarto es el diagnóstico de que, entre dominios, lo que se rompe primero es
la calibración y no la representación, junto con una receta que cuesta 25
etiquetas (sección~\ref{sec:cross-dia}).

El quinto es la constatación de que la validación entre días no detecta la fuga
de identificadores de red, sino que la premia (sección~\ref{sec:fuga}).

A ello se añade un resultado negativo que también es una aportación: ninguna de
las arquitecturas avanzadas contrastadas mejora los resultados sobre estos
datos, y dos de ellas cuestan entre 2 y 49 veces más
(sección~\ref{sec:sota-contraste}). El límite no está en la capacidad del
modelo.

\subsection*{Cómo se llegó a esos resultados}

Queda decir algo sobre el método de trabajo, porque las cinco aportaciones
tienen un origen común y no es la planificación inicial. Ninguna estaba prevista.
Todas aparecieron al desconfiar de un resultado que ya funcionaba.

El patrón se repite y esta memoria lo documenta en lugar de esconderlo. Un
primer modelo alcanzaba el 99.95\,\% de exactitud, y auditarlo contra el tráfico
que un detector vería a diario reveló una tasa de falsos positivos del 96\,\%. El
tráfico benigno de referencia estaba contaminado por el propio atacante, dos
veces, y el error no se manifestaba como un fallo sino como un número plausible.
La saturación intra-día se propuso explicar mediante la memorización del host
atacante, se diseñó el experimento capaz de tumbar esa explicación, se le añadió
un control positivo para comprobar que el experimento era sensible, y la
explicación cayó. Una regla sobre la generalización entre días se enunció con dos
observaciones, resultó falsa al añadir la tercera y volvió a cambiar al medirla
con la métrica adecuada. Y varias cifras se citaron durante semanas como macro-F1
cuando eran exactitud.

La sección~\ref{sec:lecciones} recoge esos episodios uno a uno, con lo que cada
uno enseñó. Están en la memoria porque un resultado que nadie ha intentado
tumbar no es un resultado, y porque la herramienta de verificación descrita más
abajo nació precisamente de ellos.

\section{Código y reproducibilidad}
\label{sec:repositorio}

Todo el código con el que se obtuvieron los resultados de esta memoria está
disponible en un repositorio público:

\begin{center}
\small
\url{https://github.com/inakimore/deteccion-de-trafico-maligno-mediante-tecnicas-de-aprendizaje-automatico-en-sistemas-ids-de-red}
\end{center}

Contiene el \emph{pipeline} completo, desde el saneado de la captura hasta el
etiquetado, las implementaciones de los modelos, los \emph{scripts} de cada
experimento y los informes que producen. El fichero \texttt{README.md} detalla
cómo ejecutar cada paso.

Merece señalarse una decisión de organización, porque sostiene lo que la
memoria afirma. Ninguna cifra de este documento está escrita a mano sin
respaldo: todas proceden de un módulo único, \texttt{resultados.py}, que recoge
lo que producen los experimentos, y una herramienta de verificación recorre los
capítulos y las tablas comprobando que cada número aparece allí. Esa herramienta
tiene además su propio test negativo, que le inyecta un error deliberado y
comprueba que lo detecta, para que un resultado en verde no pueda ser un
falso verde. La sección~\ref{sec:lecciones} explica de qué errores nació.

Lo que el repositorio no contiene son las capturas en bruto, que ocupan unos 150
gigabytes y se descargan del organismo que las distribuye \cite{CSECICIDS2018}.

\section{Estructura de la memoria}
\label{sec:estructura}

El capítulo~\ref{ch:estado-del-arte} sitúa el trabajo en su campo. Repasa las
familias de modelos empleadas en detección de intrusiones, indicando en cada
caso si esta memoria la ha contrastado, y dedica la segunda mitad a los
conjuntos de datos públicos y a sus defectos documentados, que es donde enlaza
con los capítulos de resultados.

El capítulo~\ref{ch:metodologia} describe cómo se pasa de una captura en bruto a
un conjunto de ejemplos etiquetados, y se detiene en la parte que suele darse
por supuesta: de dónde sale la etiqueta. Contiene la verificación de la
\emph{ground-truth} y el episodio del tráfico benigno contaminado.

El capítulo~\ref{ch:regimenes} mide. Compara las tres representaciones sobre los
seis días con un único protocolo, muestra que la comparación está saturada,
contrasta y refuta la explicación que primero se propuso para esa saturación, y
termina en los ataques que el análisis de contenido no puede ver.

El capítulo~\ref{ch:generalizacion} inventaria primero las técnicas de
aprendizaje automático empleadas, tanto las propias como las tomadas del estado
del arte para contrastarlas, y responde después a la pregunta que el capítulo
anterior deja abierta: si dentro del día no se distinguen, cuál resiste cuando
cambian las condiciones. Aplica cuatro perturbaciones de exigencia creciente,
que son la interpretabilidad, la evasión, el cambio de día y el cambio de
conjunto de datos, y cierra contrastando las arquitecturas del estado del arte.

El capítulo~\ref{ch:conclusiones} cierra los objetivos uno a uno, recoge las
lecciones metodológicas, entre ellas las veces que este trabajo se corrigió a sí
mismo, declara las limitaciones y propone las líneas de continuación.
